\section{Proofs from Section \ref{subsec:long-time}}
\label{sec:proof-long-time}

Throughout this section, the hypotheses of
Theorem~\ref{theorem:trajectory} are in force.
We first record fixed-time consequences of that theorem and prove
the divergence of exact identification times.
Recall
$\tilde{\mathbb P}_{\mathcal U}
=\tilde{\mathbb P}(\,\cdot\mid\tilde\Delta=0)$.
For $\mathcal A\in\{\mathcal L,\mathcal U\}$, let 
$n_{\mathcal L}:=N$ and $n_{\mathcal U}:=M$, and write
$X_{\mathcal A}$ and $\breve X_{\mathcal A}$ for the matrices
with rows $X_i^\top$ and $\breve X_i^\top$, respectively,
$i\in I_{\mathcal A}$. We also
write $Y_{\mathcal A}:=(Y_i)_{i\in I_{\mathcal A}}$ and let
$U_{\mathcal A}$ denote the corresponding row restriction
of the Gaussian design-noise matrix. Then
\[
    \breve X_{\mathcal A}P_d^{1/2}
    =[\,\mathbf1_{n_{\mathcal A}}\;\;X_{\mathcal A}\,].
\]
We also use the norm
$\|u\|_{P_d}:=\|P_d^{1/2}u\|$.
For every $v\in\mathbb R^{n_{\mathcal A}}$,
\(
    \|\breve X_{\mathcal A}^{\top}v\|_{P_d}
    =
    \|(\breve X_{\mathcal A}P_d^{1/2})^\top v\|
\).
Write the scalar configuration rule as
\[
    \mathfrak c_\kappa(r)
    :=
    \mathbf1_{\{r\geq\kappa_+\}}
    -
    \mathbf1_{\{r\leq\kappa_-\}},
\]
and, for integers $s,t\geq1$, set
\[
    \mathfrak c^t:=\mathfrak C(\theta^t),
    \qquad
    D^t:=\|\theta^{t+1}-\theta^t\|_{P_d^{-1}}^2,
    \qquad
    h(s,t):=\frac1M\sum_{j\in I_{\mathcal U}}
        \mathbf1_{\{\mathfrak c_j^s\neq\mathfrak c_j^t\}}.
\]
For every fixed $s,t\geq1$,
Theorem~\ref{theorem:trajectory} gives
\begin{align}
    \|\theta^t-\theta^s\|_{P_d^{-1}}^2
    &\xrightarrow{\mathbb P}
    \|\tilde\theta^t-\tilde\theta^s\|_{\mathcal H}^2,
    \label{eq:long-time-parameter-distance-limit}\\
    \frac1n\|r^t-r^s\|^2
    &\xrightarrow{\mathbb P}
    \tilde{\mathbb E}
        [|\tilde r^t-\tilde r^s|^2],
    \label{eq:long-time-score-distance-limit}\\
    h(s,t)&\xrightarrow{\mathbb P}\bar h(s,t).
    \label{eq:long-time-configuration-limit}
\end{align}
The first two limits follow from the joint empirical
$W_2$ convergence, intercept convergence, and
Corollary~\ref{corollary:wasserstein-empirical-observables}.
For the last limit, apply part~2 of
Lemma~\ref{lemma:mf-proof-empirical-stability}
to the joint array $(\Delta,r^s,r^t)$, with coefficient
$(1-\rho)^{-1}\to(1-\bar\rho)^{-1}$.
The configuration-discrepancy indicator is bounded and
is discontinuous only when either score equals a selection
threshold. These events have limiting probability zero by
Lemma~\ref{lemma:mf-proof-nondegeneracy}.

\subsection{Proof of intermediate results}


\begin{proof}[Proof of Proposition~\ref{prop:identification-time-divergence}]
For $t\geq1$, recall the sample-side history
$\mathcal H_{t+1}^{[g]}$ from the proof of
Lemma~\ref{lemma:mf-proof-nondegeneracy}.
The fresh-noise representation
\eqref{eq:mf-proof-score-fresh-noise}, established in the
proof of Lemma~\ref{lemma:mf-proof-nondegeneracy}, gives
\begin{equation}
    \tilde r^{t+1}
    =
    \tilde{\mathbb E}[ \tilde r^{t+1} \mid \mathcal{H}^{[g]}_{t+1} ]
    +\sigma\bar v_{t+1}^{[w]}z_{t+1}^{[g]},
    \qquad
    \bar v_{t+1}^{[w]}>0,
    \label{eq:long-time-score-innovation}
\end{equation}
where $z_{t+1}^{[g]}$ is independent of
$\mathcal H_{t+1}^{[g]}$.
The conditional score law therefore assigns positive probability
to each of the three selection regions.
Since $\tilde{\mathfrak c}^{\,t}$ is
$\mathcal H_{t+1}^{[g]}$-measurable,
\[
    \tilde{\mathbb P}\left(
        \tilde{\mathfrak c}^{\,t+1}
        \neq\tilde{\mathfrak c}^{\,t}
        \mid\mathcal H_{t+1}^{[g]}
    \right)>0
    \qquad\text{almost surely}.
\]
Conditioning on $\tilde\Delta=0$ gives
$\bar h(t,t+1)>0$.
For every fixed integer $T\geq1$,
\[
    \{T_{\mathrm{id}}\leq T\}
    \subseteq\{h(T,T+1)=0\}.
\]
By \eqref{eq:long-time-configuration-limit},
$h(T,T+1)$ converges in probability to the strictly positive
constant $\bar h(T,T+1)$. Hence
\[
    \mathbb P(T_{\mathrm{id}}\leq T)
    \leq\mathbb P(h(T,T+1)=0)
    \longrightarrow0.
\]
\end{proof}

\subsection{Proof of Theorem~\ref{thm:effective-process-convergence}}

Throughout this subsection, assume the hypotheses of
Theorem~\ref{thm:effective-process-convergence}.
We first establish a uniform upper bound on branch Hessians
in Lemma~\ref{lem:st-deterministic-curvature-bound}
and transfer the finite-dimensional descent inequality to the
mean-field limit in
Section~\ref{subsubsec:mean-field-potential-descent}.
We then control configuration switching in
Section~\ref{subsubsec:mean-field-configuration-switching}
and derive a two-time stability estimate.
Together with square summability of the increments, these
estimates prove parameter and score convergence in
Section~\ref{subsubsec:convergence-effective-process}.

By \eqref{eq:mean-field-long-time-step-size}, we may fix
a deterministic constant $\Lambda_*$ satisfying
\[
    \Lambda_{\mathrm{bd}}<\Lambda_*<\frac2\eta.
\]
The following lemma justifies this deterministic curvature bound.

\begin{lemma}[Deterministic curvature bound]
\label{lem:st-deterministic-curvature-bound}
For every deterministic $\Lambda_*>\Lambda_{\mathrm{bd}}$,
with $\Lambda_d$ defined in \eqref{eq:lambda-d}, we have,
as $d\to\infty$,
\begin{equation}
    \mathbb P(\Lambda_d\leq\Lambda_*)\longrightarrow1.
    \label{eq:st-deterministic-curvature-bound}
\end{equation}
\end{lemma}

\begin{proof}
Since
$\breve X_{\mathcal A}P_d^{1/2}
=[\,\mathbf1_{n_{\mathcal A}}\;\;X_{\mathcal A}\,]$,
the definition of $H_d$ gives
\begin{equation}
    \Lambda_d
    \leq\lambda L_j
       +L_\ell\left(
           1+\frac{\|X_{\mathcal L}\|_{\mathrm{op}}^2}{N}
       \right)
       +\pi L_\ell\left(
           1+\frac{\|X_{\mathcal U}\|_{\mathrm{op}}^2}{M}
       \right).
    \label{eq:long-time-curvature-upper}
\end{equation}
The Gaussian-mixture representation and
$\|Y_{\mathcal A}\|=\sqrt{n_{\mathcal A}}$ imply
\[
    \frac{\|X_{\mathcal A}\|_{\mathrm{op}}}
         {\sqrt{n_{\mathcal A}}}
    \leq
    \frac{\|\mu\|}{\sqrt d}
    +\sigma\frac{\|U_{\mathcal A}\|_{\mathrm{op}}}
                      {\sqrt{n_{\mathcal A}}}.
\]
The Gaussian operator-norm bound
\cite[Theorem~5.32 and Proposition~5.34]{vershynin2012introduction}
gives
\[
    \mathbb P\left(
        \|U_{\mathcal A}\|_{\mathrm{op}}
        >\sqrt{n_{\mathcal A}}+\sqrt d+u
    \right)
    \leq e^{-u^2/2},
    \qquad u\geq0.
\]
Together with the proportional limits and signal-norm convergence,
this yields, for every $\epsilon>0$,
\begin{equation}
    \mathbb P\left(
        \frac{\|X_{\mathcal A}\|_{\mathrm{op}}}
             {\sqrt{n_{\mathcal A}}}
        >\bar K_{\mathcal A}+\epsilon
    \right)\longrightarrow0.
    \label{eq:long-time-design-bound}
\end{equation}
Combining this estimate with
\eqref{eq:long-time-curvature-upper} and
$\Lambda_*>\Lambda_{\mathrm{bd}}$ proves the claim.
\end{proof}

\subsubsection{Mean-field potential descent}
\label{subsubsec:mean-field-potential-descent}

For $t\geq1$, define the mean-field potential by
\begin{align}
    \bar V^t
    &:=
    \frac1{\bar\rho}
        \tilde{\mathbb E}
        [\tilde\Delta\,\ell(\tilde Y,\tilde r^t)]
    +
    \frac{\pi}{1-\bar\rho}
        \tilde{\mathbb E}
        [(1-\tilde\Delta)\Psi_\kappa(\tilde r^t)]
    +
    \lambda\tilde{\mathbb E}[j(\tilde w^t)].
    \label{eq:mean-field-potential}
\end{align}

\begin{proposition}[Mean-field potential descent]
\label{prop:mean-field-potential-descent}
Set $c_*:=\eta^{-1}-\Lambda_*/2>0$. For every $t\geq1$,
\begin{equation}
    \bar V^{t+1}-\bar V^t
    \leq-c_*\|\tilde\theta^{t+1}-\tilde\theta^t\|_{\mathcal H}^2.
    \label{eq:mean-field-potential-descent}
\end{equation}
Consequently, $(\bar V^t)_{t\geq1}$ decreases to a finite limit
$\bar V^\infty$, and
\begin{equation}
    \sum_{t=1}^{\infty}
        \|\tilde\theta^{t+1}-\tilde\theta^t\|_{\mathcal H}^2
    \leq\frac{\bar V^1-\bar V^\infty}{c_*}<\infty.
    \label{eq:mean-field-increment-summability}
\end{equation}
\end{proposition}

\begin{proof}
The integrands defining the potential are continuous and have
at most quadratic growth. Theorem~\ref{theorem:trajectory} and
Corollary~\ref{corollary:wasserstein-empirical-observables}
therefore give, for every fixed $t\geq1$,
\[
    V(\theta^t)\xrightarrow{\mathbb P}\bar V^t.
\]
Moreover, \eqref{eq:long-time-parameter-distance-limit} gives $D^t\xrightarrow{\mathbb P}
    \|\tilde\theta^{t+1}-\tilde\theta^t\|_{\mathcal H}^2$. 
On $\{\Lambda_d\leq\Lambda_*\}$, the finite-dimensional
descent inequality yields
\[
    V(\theta^{t+1})-V(\theta^t)\leq-c_*D^t.
\]
This event has probability tending to one, so passing to the
deterministic limits proves
\eqref{eq:mean-field-potential-descent}.
Recall that
\[
    -C_\kappa\leq\Psi_\kappa\leq0,
    \qquad
    C_\kappa:=
    \max\{\ell(-1,\kappa_-),\ell(+1,\kappa_+)\}.
\]
Nonnegativity of the supervised loss and penalty gives
$\bar V^t\geq-\pi C_\kappa$.
Thus the decreasing potential has a finite limit.
Summing \eqref{eq:mean-field-potential-descent} proves
\eqref{eq:mean-field-increment-summability}.
\end{proof}

Square summability controls consecutive iterates but does not
by itself imply that the trajectory is Cauchy. We next control
changes of configuration and then compare arbitrary pairs of
iterates. For these finite-dimensional estimates, fix
\begin{align}
    0<a<\min\{\bar p_{\mathcal L},1-\bar p_{\mathcal L}\},
    \quad K_{\mathcal L}>\bar K_{\mathcal L}
    , \quad K_{\mathcal U}>\bar K_{\mathcal U},
    \quad K_{\mathrm{aug}}:=\sqrt{1+K_{\mathcal U}^2}, \label{eq:a-k-k-aug}
\end{align}
with $\bar K_{\mathcal{L}}$ and $\bar K_{\mathcal U}$ given in \eqref{eq:s-mu-k-bar},
and define
\begin{equation}
\begin{split}
    \mathcal E_d:=\Bigl\{
        p_{\mathcal L}\in[a,1-a],\
        \|X_{\mathcal L}\|_{\mathrm{op}}\leq K_{\mathcal L}\sqrt N, \
        \|X_{\mathcal U}\|_{\mathrm{op}}\leq K_{\mathcal U}\sqrt M,\
        \Lambda_d\leq\Lambda_*,\
        V(\theta^1)\leq\bar V^1+1
    \Bigr\}.
\end{split}
    \label{eq:long-time-good-event}
\end{equation}
The preceding bounds imply $\mathbb P(\mathcal E_d)\to1$.
On this event,
\begin{equation}
    \|\breve X_{\mathcal U}P_d^{1/2}\|_{\mathrm{op}}
    \leq K_{\mathrm{aug}}\sqrt M,
    \qquad
    0\preceq P_d^{1/2}\nabla^2L_{\mathfrak c}(\theta)P_d^{1/2}
       \preceq\Lambda_*I_{d+1}
    \label{eq:long-time-uniform-upper-bounds}
\end{equation}
for every $\theta$ and every configuration $\mathfrak c$.

\subsubsection{Control of configuration switching}
\label{subsubsec:mean-field-configuration-switching}

A change of active pseudo-label produces a nonzero jump in the
scalar residual. The next argument shows that the Gaussian design
prevents the aggregate effect of these jumps from becoming too
small relative to the fraction of observations that switch.
Define the $L_\ell$-Lipschitz functions
\begin{align}
    f_{-1}(r):=\ell'(-1,r),\qquad f_0(r):=0,\qquad
    f_{+1}(r):=\ell'(+1,r). \label{eq:def-f}
\end{align}

\begin{proposition}[Control of configuration switching]
\label{prop:mean-field-configuration-switching}
There exists a deterministic $C_{\mathrm{sw}}\in(0,\infty)$,
independent of iteration time, such that
\begin{equation}
    \bar h(t,t+1)
    \leq C_{\mathrm{sw}}\left(
        \|\tilde\theta^{t+1}-\tilde\theta^t\|_{\mathcal H}^2
        +\|\tilde\theta^{t+2}-\tilde\theta^{t+1}\|_{\mathcal H}^2
    \right),\qquad t\geq1.
    \label{eq:effective-switching-control}
\end{equation}
Consequently, the configuration switching probabilities are summable
\begin{align}
    \sum_{t=1}^{\infty}\bar h(t,t+1)
    &\leq2C_{\mathrm{sw}}\sum_{t=1}^{\infty}
        \|\tilde\theta^{t+1}-\tilde\theta^t\|_{\mathcal H}^2
      <\infty,
    \label{eq:effective-switching-summability}
\end{align}
and we get
\begin{align}
    \lim_{T\to\infty}\sup_{s,t\geq T}\bar h(s,t)&=0.
    \label{eq:effective-configuration-cauchy}
\end{align}
\end{proposition}

\begin{proof}
Define the threshold jump magnitudes
\[
    \varsigma_-:=\ell'(-1,\kappa_-)>0,\qquad
    \varsigma_+:=-\ell'(+1,\kappa_+)>0,\qquad
    \varsigma_\kappa:=\min\{\varsigma_-,\varsigma_+\}>0,
\]
and the finite set
$\mathcal A_\kappa:=\{0,\pm \varsigma_-,\pm \varsigma_+,\pm(\varsigma_-+\varsigma_+)\}$.
We prove the result in three steps.

\begin{itemize}[wide]
    \item \emph{Step 1: A uniform Gaussian lower bound.}
There exist a deterministic $\gamma>0$ and events
$\mathcal R_d$ with $\mathbb P(\mathcal R_d)\to1$ such that,
for all sufficiently large $d$, on $\mathcal R_d$,
\begin{equation}
    \|\breve X_{\mathcal U}^{\top}v\|_{P_d}
    \geq\gamma\sqrt M\,\|v\|,
    \qquad v\in\mathcal A_\kappa^M.
    \label{eq:configuration-switching-gaussian-bound}
\end{equation}
To see this, for $d\geq2$ choose a deterministic
$Q\in\mathbb R^{d\times(d-1)}$ with
$Q^\top Q=I_{d-1}$ and $Q^\top\mu=0$. Then
\[
    Q^\top X_{\mathcal U}^{\top}
    =\sigma Q^\top U_{\mathcal U}^{\top},
\]
and $Q^\top U_{\mathcal U}^{\top}$ has independent standard
Gaussian entries. Thus, for each deterministic nonzero $v$,
\[
    \frac{Q^\top U_{\mathcal U}^{\top}v}{\|v\|}
    \sim\mathcal N(0,I_{d-1}).
\]
For $Z\sim\mathcal N(0,I_{d-1})$, exponential Markov inequality
gives, for every $\varepsilon>0$,
\[
    \mathbb P(\|Z\|^2\leq\varepsilon^2(d-1))
    \leq e^{(d-1)/2}
       \mathbb E[e^{-\|Z\|^2/(2\varepsilon^2)}]
    =e^{(d-1)/2}(1+\varepsilon^{-2})^{-(d-1)/2}
      \leq(\sqrt e\,\varepsilon)^{d-1}.
\]
A union bound over at most $7^M$ vectors therefore yields
\[
\begin{split}
    \mathbb P\Bigl(
        \exists v\in\mathcal A_\kappa^M\setminus\{0\}:\,
        \|Q^\top U_{\mathcal U}^{\top}v\|
        <\varepsilon\sqrt{d-1}\,\|v\|
    \Bigr)
    \leq7^M(\sqrt e\,\varepsilon)^{d-1}.
\end{split}
\]
Since $M/(d-1)\to(1-\bar\rho)\bar\delta\in(0,\infty)$,
we may fix $\varepsilon>0$ small enough that the right-hand
side is at most $e^{-(d-1)}$ for all sufficiently large $d$.
Define
\[
    \mathcal R_d
    :=
    \bigcap_{v\in\mathcal A_\kappa^M\setminus\{0\}}
    \left\{
        \|Q^\top U_{\mathcal U}^{\top}v\|
        \geq\varepsilon\sqrt{d-1}\,\|v\|
    \right\}.
\]
The preceding union bound gives, for all sufficiently large $d$,
\[
    \mathbb P(\mathcal R_d^c)
    \leq 7^M(\sqrt e\,\varepsilon)^{d-1}
    \leq e^{-(d-1)}.
\]
Hence $\mathbb P(\mathcal R_d)\to1$.
Finally, we may also fix $\gamma>0$ such that
$\sigma\varepsilon\sqrt{d-1}\geq\gamma\sqrt M$ for all
sufficiently large $d$, thus on $\mathcal{R}_d$, we have

\[
    \|P_d^{1/2}\breve X_{\mathcal U}^{\top}v\|^2
    =|\mathbf1_M^\top v|^2+\|X_{\mathcal U}^{\top}v\|^2
    \geq\sigma^2\|Q^\top U_{\mathcal U}^{\top}v\|^2,
\]
which proves \eqref{eq:configuration-switching-gaussian-bound}.
The bound holds simultaneously for every vector in the finite
set, so it also applies to vectors chosen from that set using
the observed data.

\item \emph{Step 2: Bounding switches by consecutive updates.}
Work on $\mathcal E_d\cap\mathcal R_d$ and fix $t\geq1$.
At the common score $r_j^{t+1}$, let
\[
    e_j^t:=f_{\mathfrak c_j^{t+1}}(r_j^{t+1})
             -f_{\mathfrak c_j^t}(r_j^{t+1}),
    \qquad j\in I_{\mathcal U}.
\]
Approximate these residual changes by $v^t\in\mathcal A_\kappa^M$,
where
\[
    v_j^t:=
    \begin{cases}
        0,&\mathfrak c_j^{t+1}=\mathfrak c_j^t,\\[2pt]
        \varsigma_-(\mathfrak c_j^t-\mathfrak c_j^{t+1}),
           &\{\mathfrak c_j^t,\mathfrak c_j^{t+1}\}=\{-1,0\},\\[2pt]
        \varsigma_+(\mathfrak c_j^t-\mathfrak c_j^{t+1}),
           &\{\mathfrak c_j^t,\mathfrak c_j^{t+1}\}=\{0,+1\},\\[2pt]
        \frac{\varsigma_-+\varsigma_+}{2}(\mathfrak c_j^t-\mathfrak c_j^{t+1}),
           &\{\mathfrak c_j^t,\mathfrak c_j^{t+1}\}=\{-1,+1\}.
    \end{cases}
\]
Every changed coordinate has magnitude at least $\varsigma_\kappa$,
so
\begin{equation}
    \|v^t\|^2\geq \varsigma_\kappa^2 M h(t,t+1).
    \label{eq:configuration-switching-support}
\end{equation}
Moreover,
\begin{equation}
    |e_j^t-v_j^t|\leq 2L_\ell|r_j^{t+1}-r_j^t|.
    \label{eq:configuration-switching-jump-error}
\end{equation}
For a transition involving $0$, the relevant threshold lies between
the two scores, and the estimate follows from the Lipschitz
continuity of the corresponding loss derivative.
For a transition from $-1$ to $+1$, write $r'=r_j^{t+1}$.
Then
\[
\begin{aligned}
    e_j^t-v_j^t
    &=
    [\ell'(+1,r')-\ell'(+1,\kappa_+)]
    -[\ell'(-1,r')-\ell'(-1,\kappa_-)].
\end{aligned}
\]
Both thresholds lie between $r_j^t$ and $r'$; hence each bracket
is bounded in absolute value by $L_\ell|r'-r_j^t|$.
The reverse transition is treated identically, and the estimate
also holds at threshold endpoints.

By the design-operator bound,
\begin{equation}
    \frac{\|e^t-v^t\|}{\sqrt M}
    \leq 2L_\ell
        \frac{\|r_{\mathcal U}^{t+1}-r_{\mathcal U}^t\|}{\sqrt M}
    \leq 2L_\ell K_{\mathrm{aug}}\sqrt{D^t},
    \label{eq:configuration-switching-vector-error}
\end{equation}
where $r_{\mathcal U}^t:=(r_j^t)_{j\in I_{\mathcal U}}$.
The same Hessian bound gives, for every configuration,
\[
    \|\nabla L_{\mathfrak c}(\theta)
          -\nabla L_{\mathfrak c}(\vartheta)\|_{P_d} \leq \int_0^1  \| P_d^{1/2} \nabla^2 L_{\mathfrak c}(\vartheta + s(\theta - \vartheta)) P_d^{1/2} P_d^{-1/2} (\theta - \vartheta) \| \, ds
    \leq\Lambda_*\|\theta-\vartheta\|_{P_d^{-1}}.
\]
The supervised and penalty gradients cancel when two branch
objectives are compared at the same parameter. Hence
\[
    \frac\pi M\breve X_{\mathcal U}^{\top}e^t
    =\nabla L_{\mathfrak c^{t+1}}(\theta^{t+1})
       -\nabla L_{\mathfrak c^t}(\theta^{t+1}).
\]
Using
$\|\nabla L_{\mathfrak c^t}(\theta^t)\|_{P_d}
=\sqrt{D^t}/\eta$ and adding and subtracting
$\nabla L_{\mathfrak c^t}(\theta^t)$ yields
\begin{equation}
    \left\|\frac\pi M\breve X_{\mathcal U}^{\top}e^t\right\|_{P_d}
    \leq\left(\Lambda_*+\frac1\eta\right)\sqrt{D^t}
         +\frac1\eta\sqrt{D^{t+1}}.
    \label{eq:configuration-switching-gradient-upper}
\end{equation}
On the other hand,
\eqref{eq:configuration-switching-gaussian-bound} and
\eqref{eq:configuration-switching-support} imply
\[
    \left\|\frac\pi M\breve X_{\mathcal U}^{\top}v^t\right\|_{P_d}
    \geq\pi\gamma \varsigma_\kappa\sqrt{h(t,t+1)},
\]
while \eqref{eq:configuration-switching-vector-error} gives
\[
    \left\|\frac\pi M\breve X_{\mathcal U}^{\top}(v^t-e^t)\right\|_{P_d}
    \leq 2\pi L_\ell K_{\mathrm{aug}}^2\sqrt{D^t}.
\]
Combining the last three estimates and squaring gives
\begin{equation}
    h(t,t+1)\leq C_{\mathrm{sw}}(D^t+D^{t+1}),
    \qquad
    C_{\mathrm{sw}}
    :=\frac{2}{\pi^2\gamma^2\varsigma_\kappa^2}
       \left(\Lambda_*+\frac1\eta
                         +2\pi L_\ell K_{\mathrm{aug}}^2\right)^2.
    \label{eq:configuration-switching-bound}
\end{equation}
All constants are deterministic and independent of $d$ and $t$.
The event $\mathcal E_d\cap\mathcal R_d$ does not depend on
time, and the estimate holds there for every $t\geq1$.

\item \emph{Step 3: Passage to the effective process.}
For fixed $t$, apply Theorem~\ref{theorem:trajectory} on the
horizon $t+2$. The convergence of $D^t$ and $D^{t+1}$, together
with \eqref{eq:long-time-configuration-limit}, allows
\eqref{eq:configuration-switching-bound} to pass to deterministic
limits, since $\mathbb P(\mathcal E_d\cap\mathcal R_d)\to1$.
This proves \eqref{eq:effective-switching-control}. Summing over
$t$ and applying \eqref{eq:mean-field-increment-summability}
gives \eqref{eq:effective-switching-summability}.
For integers $t>s$,
\[
    \{\tilde{\mathfrak c}^{\,s}\neq\tilde{\mathfrak c}^{\,t}\}
    \subseteq\bigcup_{k=s}^{t-1}
        \{\tilde{\mathfrak c}^{\,k+1}\neq\tilde{\mathfrak c}^{\,k}\}.
\]
Therefore
\[
    \sup_{s,t\geq T}\bar h(s,t)
    \leq\sum_{k=T}^{\infty}\bar h(k,k+1)\longrightarrow0,
\]
proving \eqref{eq:effective-configuration-cauchy}.

\end{itemize}

\end{proof}

\subsubsection{Convergence of the effective process}
\label{subsubsec:convergence-effective-process}

To compare parameters at different times, we need a quantitative
version of branch convexity. The strongly convex penalty controls weight
directions. The labeled loss supplies the remaining curvature in
the intercept direction once the trajectory is uniformly localized.

\begin{lemma}[Uniform localization and branch curvature]
\label{lem:uniform-branch-curvature}
Set
\[
    C:=\bar V^1+\pi C_\kappa+1,
    \qquad
    \mathcal K_d:=\{\theta \in \mathbb{R}^{d+1}:L_{\mathrm{sup}}(\theta)\leq C\},
    \qquad
    L_{\mathrm{sup}}(\theta)
    :=L_{\mathcal L}(\theta)+\frac{\lambda}{d}J(w).
\]
There exists a deterministic $m_*>0$, independent of $d$
and the configuration, such that, on $\mathcal E_d$,
the convex set $\mathcal K_d$ contains every $\theta^t$,
$t\geq1$, and
\begin{equation}
    \nabla^2L_{\mathfrak c}(\theta)\succeq m_*P_d^{-1},
    \qquad\theta\in\mathcal K_d,
    \qquad\mathfrak c\in\{-1,0,+1\}^{I_{\mathcal U}}.
    \label{eq:long-time-uniform-lower-curvature}
\end{equation}
\end{lemma}

\begin{proof}
Work on $\mathcal E_d$. Since $\Psi_\kappa\geq-C_\kappa$ from \eqref{eq:c-kappa}, the
finite-dimensional potential descent Lemma~\ref{lem:st-global-descent-boundedness} gives
\[
    L_{\mathrm{sup}}(\theta^t)
    \leq V(\theta^t)+\pi C_\kappa
    \leq V(\theta^1)+\pi C_\kappa\leq C,
    \qquad t\geq1.
\]
The set $\mathcal K_d$ is convex because $L_{\mathrm{sup}}$
is convex. Since $j\geq0$ and $j(0)=0$, differentiability gives
$j'(0)=0$. The bound $j''\geq m_j$ therefore implies
$j(u)\geq m_j u^2/2$.
Thus, for $\theta=(b,w)\in\mathcal K_d$, nonnegativity of
the loss $L_{\mathcal L}$ gives
\[
    \lambda \frac{m_j}{2} \frac{\| w \|^2}{d} \le \frac{\lambda}{d}J(w) \le L_{\mathrm{sup}}(\theta) \le C, 
\]
hence
\[
    \frac{\|w\|}{\sqrt d}
    \le   W:=\sqrt{\frac{2C}{\lambda m_j}}.
\]
To bound the intercept, define
\[
\begin{aligned}
    a_-&:=\ell'(-1,0)>0,\qquad
    a_+:=-\ell'(+1,0)>0,\\
    a_{\min}&:=\min\{a_-,a_+\},\qquad
    a_{\max}:=\max\{a_-,a_+\}.
\end{aligned}
\]
Convexity gives the tangent-line bounds
\[
\begin{aligned}
    \ell(-1,r)
    &\geq \ell(-1,0)+a_-r
     \geq a_-r,&
    \ell(+1,r)
    &\geq \ell(+1,0)-a_+r
     \geq -a_+r,
\end{aligned}
\]
where the second inequalities use nonnegativity of the loss.
Let
\[
    I_\pm:=\{i\in I_{\mathcal L}:Y_i=\pm1\},
    \qquad
    z_i:=\frac{X_i^\top w}{\sqrt d}.
\]
On $\mathcal E_d$, both labeled classes have proportion at least
$a$: $\#I_- / N, \ \# I_+/ N \geq a$ where $a$ is fixed in \eqref{eq:a-k-k-aug}.
Since every loss term is nonnegative, summing over either class
alone gives a lower bound on $L_{\mathcal L}$.
If $b\geq0$, use the negatively labeled observations to obtain
\begin{align}
    L_{\mathcal L}(\theta)
    \geq \frac1N\sum_{i\in I_-}\ell(-1,b+z_i)
    \geq a_-\frac{|I_-|}{N}b
           -\frac{a_-}{N}\sum_{i\in I_-}|z_i|
    \geq a\,a_{\min}|b|
           -\frac{a_{\max}}{N}
              \sum_{i\in I_{\mathcal L}}|z_i|.
    \label{eq:labeled-loss}
\end{align}
If $b<0$, we get the same bound using $I_+$ instead of $I_-$.
By Cauchy--Schwarz and the design bound on $\mathcal E_d$,
\[
    \frac1N\sum_{i\in I_{\mathcal L}}|z_i|
    \leq \frac{\|X_{\mathcal L}w\|}{\sqrt{Nd}}
    \leq
    \frac{\|X_{\mathcal L}\|_{\mathrm{op}}}{\sqrt N}
    \frac{\|w\|}{\sqrt d}
    \leq K_{\mathcal L}\frac{\|w\|}{\sqrt d}.
\]
Thus, \eqref{eq:labeled-loss} becomes
\[
    L_{\mathcal L}(\theta)
    \geq a\,a_{\min}|b|
        -a_{\max}K_{\mathcal L}\frac{\|w\|}{\sqrt d}.
\]
For $\theta\in\mathcal K_d$, we have
$L_{\mathcal L}(\theta)\leq L_{\mathrm{sup}}(\theta)\leq C$
and $\|w\|/\sqrt d\leq W$. Consequently,
\[
    |b|
    \leq B:=
    \frac{C+a_{\max}K_{\mathcal L}W}{a\,a_{\min}}.
\]
Finally, using $r_i(\theta)=b+z_i$ and
$(x+y)^2\leq2x^2+2y^2$, we obtain
\[
    \frac1N\sum_{i\in I_{\mathcal L}}r_i(\theta)^2
    \leq 2b^2+\frac{2\|X_{\mathcal L}w\|^2}{Nd}
    \leq 2B^2+2K_{\mathcal L}^2W^2
    =:Q_*.
\]
Set $R:=\sqrt{2Q_*}$. The preceding score bound gives
\[
    \frac1N
    \#\{i\in I_{\mathcal L}:|r_i(\theta)|>R\}
    \leq
    \frac1{NR^2}\sum_{i\in I_{\mathcal L}}r_i(\theta)^2
    \leq\frac12.
\]
Thus at least half of the labeled scores belong to $[-R,R]$.
By continuity and strict positivity of the loss curvature,
\[
    \alpha_R:=
    \min_{\substack{y\in\{-1,+1\}\\|r|\leq R}}
        \ell''(y,r)>0.
\]
Keeping the curvature contributions from these bounded scores
and using nonnegativity of the remaining contributions yields
\begin{equation}
    \frac1N\sum_{i\in I_{\mathcal L}}
        \ell''(Y_i,r_i(\theta))
    \geq a_0:=\frac{\alpha_R}{2}>0.
    \label{eq:long-time-intercept-curvature}
\end{equation}
For a direction $(h,u)\in\mathbb R\times\mathbb R^d$, define
the directional curvature of the labeled loss by
\[
    Q_{\mathcal L}(\theta;h,u)
    :=
    \frac1N\sum_{i\in I_{\mathcal L}}
        \ell''(Y_i,r_i(\theta))
        \left(h+\frac{X_i^\top u}{\sqrt d}\right)^2.
\]
Although \eqref{eq:long-time-intercept-curvature} controls
pure intercept directions, the two terms inside the square
may cancel for a mixed direction. To control this cancellation,
write
$h=
    \left(h+{X_i^\top u}/{\sqrt d}\right)
    -{X_i^\top u}/{\sqrt d}$.
Applying the inequality $(x-y)^2\leq2x^2+2y^2$, multiplying by
$\ell''(Y_i,r_i(\theta))$, and averaging gives
\begin{align*}
    a_0h^2
    &\leq
    \frac1N\sum_{i\in I_{\mathcal L}}
        \ell''(Y_i,r_i(\theta))h^2
    \leq
    2Q_{\mathcal L}(\theta;h,u)
    +\frac2{Nd}\sum_{i\in I_{\mathcal L}}
        \ell''(Y_i,r_i(\theta))(X_i^\top u)^2\\
    &\leq
    2Q_{\mathcal L}(\theta;h,u)
    +\frac{2L_\ell}{Nd}\|X_{\mathcal L}u\|^2
    \leq
    2Q_{\mathcal L}(\theta;h,u)
    +2L_\ell K_{\mathcal L}^2\frac{\|u\|^2}{d}.
\end{align*}
Dividing by $a_0$ and adding $\|u\|^2/d$ yields
\[
    h^2+\frac{\|u\|^2}{d}
    \leq
    \frac2{a_0}Q_{\mathcal L}(\theta;h,u)
    +C_0\frac{\|u\|^2}{d},
    \qquad
    C_0:=1+\frac{2L_\ell K_{\mathcal L}^2}{a_0}.
\]
The unlabeled contribution to the branch Hessian is positive
semidefinite, while $j''\geq m_j$ gives
\[
\begin{aligned}
    \begin{pmatrix}h\\u\end{pmatrix}^{\!\top}
    \nabla^2L_{\mathfrak c}(\theta)
    \begin{pmatrix}h\\u\end{pmatrix}
    \geq
    Q_{\mathcal L}(\theta;h,u)
    +\frac{\lambda}{d}\sum_{k=1}^d j''(w_k)u_k^2
    \geq
    Q_{\mathcal L}(\theta;h,u)
    +\lambda m_j\frac{\|u\|^2}{d}.
\end{aligned}
\]
Choose
\[
    m_*:=
    \min\left\{
        \frac{a_0}{2},
        \frac{\lambda m_j}{C_0}
    \right\}
    =
    \min\left\{
        \frac{a_0}{2},
        \frac{\lambda m_j}
             {1+2L_\ell K_{\mathcal L}^2/a_0}
    \right\}>0.
\]
Since $2m_*/a_0\leq1$ and $m_*C_0\leq\lambda m_j$,
\[
\begin{aligned}
    m_*\left(h^2+\frac{\|u\|^2}{d}\right)
    \leq
    \frac{2m_*}{a_0}Q_{\mathcal L}(\theta;h,u)
    +m_*C_0\frac{\|u\|^2}{d}
    \leq
    Q_{\mathcal L}(\theta;h,u)
    +\lambda m_j\frac{\|u\|^2}{d}
    \leq
    \begin{pmatrix}h\\u\end{pmatrix}^{\!\top}
    \nabla^2L_{\mathfrak c}(\theta)
    \begin{pmatrix}h\\u\end{pmatrix}.
\end{aligned}
\]
This proves
\eqref{eq:long-time-uniform-lower-curvature}
on $\mathcal E_d$, uniformly over
$\theta\in\mathcal K_d$ and all configurations $\mathfrak c$.
\end{proof}

\begin{lemma}[Two-time stability]
\label{lem:effective-two-time-stability}
There exists a deterministic $A\in(0,\infty)$ such that,
for every $s,t\geq1$,
\begin{align}
    \|\tilde\theta^t-\tilde\theta^s\|_{\mathcal H}
    &\leq A\Bigl(
        \|\tilde\theta^{s+1}-\tilde\theta^s\|_{\mathcal H}
        +\|\tilde\theta^{t+1}-\tilde\theta^t\|_{\mathcal H}
        +\sqrt{\bar h(s,t)}\Bigr),
    \label{eq:effective-two-time-stability-main}\\
    \|\tilde r^t-\tilde r^s\|_{L^2(\tilde{\mathbb P})}
    &\leq A\|\tilde\theta^t-\tilde\theta^s\|_{\mathcal H}.
    \label{eq:effective-score-stability-main}
\end{align}
\end{lemma}

\begin{proof}
We first establish finite-dimensional estimates on
$\mathcal E_d$, with constants independent of the dimension
and iteration times, and then pass to the trajectory limit.

\begin{itemize}[wide]
    \item \emph{Step 1: Comparing branch gradients at a common parameter.}
Recall the functions $f_c$ from \eqref{eq:def-f}.
Since $f_0=0$ and $|\ell'(y,r)|\leq B_\ell$ from Assumption~\ref{assumption:long-time-loss-penalty},
\[
    |f_c(r)-f_{c'}(r)|
    \leq2B_\ell\mathbf1_{\{c\neq c'\}},
    \qquad
    c,c'\in\{-1,0,+1\},\quad r\in\mathbb R.
\]
For a common parameter $\theta$, define the residual-difference
vector $e^{s,t}(\theta)\in\mathbb R^M$ by
\[
    e_j^{s,t}(\theta)
    :=
    f_{\mathfrak c_j^t}(r_j(\theta))
    -
    f_{\mathfrak c_j^s}(r_j(\theta)),
    \qquad j\in I_{\mathcal U}.
\]
Only observations whose configurations differ can contribute.
Consequently,
\[
\begin{aligned}
    \|e^{s,t}(\theta)\|^2
    \leq
    4B_\ell^2
    \sum_{j\in I_{\mathcal U}}
        \mathbf1_{\{\mathfrak c_j^t\neq\mathfrak c_j^s\}}
    =4B_\ell^2Mh(s,t).
\end{aligned}
\]
The supervised and penalty terms are identical in all branch
objectives. They therefore cancel when the branch gradients
are evaluated at the same parameter:
\[
    \nabla L_{\mathfrak c^t}(\theta)
    -
    \nabla L_{\mathfrak c^s}(\theta)
    =
    \frac{\pi}{M}
    \breve X_{\mathcal U}^{\top}e^{s,t}(\theta).
\]
Using the design bound on $\mathcal E_d$ gives
\[
\begin{aligned}
    \|\nabla L_{\mathfrak c^t}(\theta)
          -\nabla L_{\mathfrak c^s}(\theta)\|_{P_d}
    =
    \frac{\pi}{M}
    \|P_d^{1/2}\breve X_{\mathcal U}^{\top}
        e^{s,t}(\theta)\|
    \leq
    \frac{\pi}{M}
    \|\breve X_{\mathcal U}P_d^{1/2}\|_{\mathrm{op}}
    \|e^{s,t}(\theta)\|
    \leq
    2\pi B_\ell K_{\mathrm{aug}}\sqrt{h(s,t)}.
\end{aligned}
\]
Thus, for every $\theta$,
\begin{equation}
    \|\nabla L_{\mathfrak c^t}(\theta)
          -\nabla L_{\mathfrak c^s}(\theta)\|_{P_d}
    \leq
    2\pi B_\ell K_{\mathrm{aug}}\sqrt{h(s,t)}.
    \label{eq:long-time-gradient-perturbation}
\end{equation}

\item \emph{Step 2: Controlling parameter differences.}
Both $\theta^s$ and $\theta^t$ belong to the convex set
$\mathcal K_d$. Their joining segment therefore also belongs
to $\mathcal K_d$, where
Lemma~\ref{lem:uniform-branch-curvature} gives
\[
    \nabla^2L_{\mathfrak c^t}(\theta)
    \succeq m_*P_d^{-1}.
\]
Set $\Delta\theta:=\theta^t-\theta^s$.
Integrating the Hessian along the segment yields
\[
\begin{aligned}
    \left\langle
        \nabla L_{\mathfrak c^t}(\theta^t)
        -\nabla L_{\mathfrak c^t}(\theta^s),
        \Delta\theta
    \right\rangle
    =
    \int_0^1
        \Delta\theta^\top
        \nabla^2L_{\mathfrak c^t}
            (\theta^s+\tau\Delta\theta)
        \Delta\theta\,d\tau
    \geq
    m_*\|\Delta\theta\|_{P_d^{-1}}^2.
\end{aligned}
\]
To bound the gradient difference, add and subtract
$\nabla L_{\mathfrak c^s}(\theta^s)$:
\[
\begin{aligned}
    \|\nabla L_{\mathfrak c^t}(\theta^t)
          -\nabla L_{\mathfrak c^t}(\theta^s)\|_{P_d}
    \leq
    \|\nabla L_{\mathfrak c^t}(\theta^t)\|_{P_d}
    +\|\nabla L_{\mathfrak c^s}(\theta^s)\|_{P_d}
    +
    \|\nabla L_{\mathfrak c^s}(\theta^s)
          -\nabla L_{\mathfrak c^t}(\theta^s)\|_{P_d}.
\end{aligned}
\]
The update identity gives, for every $k\geq1$,
\[
    \nabla L_{\mathfrak c^k}(\theta^k)
    =
    -\frac1\eta P_d^{-1}(\theta^{k+1}-\theta^k),
\]
and hence
\[
    \|\nabla L_{\mathfrak c^k}(\theta^k)\|_{P_d}
    =
    \frac1\eta
    \|\theta^{k+1}-\theta^k\|_{P_d^{-1}}
    =
    \frac{\sqrt{D^k}}{\eta}.
\]
Combining these identities with
\eqref{eq:long-time-gradient-perturbation} yields
\[
\begin{aligned}
    \|\nabla L_{\mathfrak c^t}(\theta^t)
          -\nabla L_{\mathfrak c^t}(\theta^s)\|_{P_d}
    \leq
    \frac{\sqrt{D^t}+\sqrt{D^s}}{\eta}
    +2\pi B_\ell K_{\mathrm{aug}}\sqrt{h(s,t)}.
\end{aligned}
\]
Cauchy--Schwarz in the dual norms gives
\[
    |\langle g,v\rangle|
    =
    |\langle P_d^{1/2}g,P_d^{-1/2}v\rangle|
    \leq\|g\|_{P_d}\|v\|_{P_d^{-1}}.
\]
Applying this inequality to the integrated curvature estimate,
we obtain
\[
\begin{aligned}
    m_*\|\Delta\theta\|_{P_d^{-1}}^2
    &\leq
    \left(
        \frac{\sqrt{D^t}+\sqrt{D^s}}{\eta}
        +2\pi B_\ell K_{\mathrm{aug}}\sqrt{h(s,t)}
    \right)
    \|\Delta\theta\|_{P_d^{-1}}.
\end{aligned}
\]
If $\Delta\theta=0$, the desired bound is immediate.
Otherwise, division by
$m_*\|\Delta\theta\|_{P_d^{-1}}$ gives
\begin{equation}
    \|\theta^t-\theta^s\|_{P_d^{-1}}
    \leq
    \frac{\sqrt{D^t}+\sqrt{D^s}}{\eta m_*}
    +\frac{2\pi B_\ell K_{\mathrm{aug}}}{m_*}
        \sqrt{h(s,t)}.
    \label{eq:long-time-finite-two-time-stability}
\end{equation}

\item \emph{Step 3: Controlling score differences.}
For either population
$\mathcal A\in\{\mathcal L,\mathcal U\}$,
\[
    r_{\mathcal A}^t-r_{\mathcal A}^s
    =
    \breve X_{\mathcal A}(\theta^t-\theta^s)
    =
    \breve X_{\mathcal A}P_d^{1/2}
    P_d^{-1/2}(\theta^t-\theta^s).
\]
The design bounds on $\mathcal E_d$ imply
\[
    \|\breve X_{\mathcal A}P_d^{1/2}\|_{\mathrm{op}}^2
    \leq
    n_{\mathcal A}+\|X_{\mathcal A}\|_{\mathrm{op}}^2
    \leq n_{\mathcal A}(1+K_{\mathcal A}^2).
\]
Consequently,
\[
    \frac{\|r_{\mathcal A}^t-r_{\mathcal A}^s\|^2}
         {n_{\mathcal A}}
    \leq
    (1+K_{\mathcal A}^2)
    \|\theta^t-\theta^s\|_{P_d^{-1}}^2.
\]
Define
\[
    K_{\mathrm{sc}}
    :=
    \sqrt{1+\max\{K_{\mathcal L}^2,K_{\mathcal U}^2\}}.
\]
Combining the two populations gives
\[
\begin{aligned}
    \frac1n\|r^t-r^s\|^2
    =
    \rho\,\frac{\|r_{\mathcal L}^t-r_{\mathcal L}^s\|^2}{N}
    +(1-\rho)\,
        \frac{\|r_{\mathcal U}^t-r_{\mathcal U}^s\|^2}{M}
    \leq
    K_{\mathrm{sc}}^2
    \|\theta^t-\theta^s\|_{P_d^{-1}}^2.
\end{aligned}
\]
Taking square roots yields
\begin{equation}
    \frac{\|r^t-r^s\|}{\sqrt n}
    \leq
    K_{\mathrm{sc}}
    \|\theta^t-\theta^s\|_{P_d^{-1}}.
    \label{eq:long-time-finite-score-stability}
\end{equation}

\item \emph{Step 4: Passage to the effective process.}
Fix $s,t\geq1$.
The limits \eqref{eq:long-time-parameter-distance-limit},
\eqref{eq:long-time-score-distance-limit}, and
\eqref{eq:long-time-configuration-limit}
apply jointly at the finitely many time pairs needed below,
including $(s,s+1)$ and $(t,t+1)$.
By continuity of the square-root function, the corresponding
unsquared distances also converge in probability.

The finite-dimensional estimates
\eqref{eq:long-time-finite-two-time-stability} and
\eqref{eq:long-time-finite-score-stability}
hold on $\mathcal E_d$, whose probability tends to one.
Since all quantities in these estimates converge in probability
to deterministic limits, the inequalities pass to those limits.
Thus
\[
\begin{aligned}
    \|\tilde\theta^t-\tilde\theta^s\|_{\mathcal H}
    \leq
    \frac{
        \|\tilde\theta^{t+1}-\tilde\theta^t\|_{\mathcal H}
        +\|\tilde\theta^{s+1}-\tilde\theta^s\|_{\mathcal H}
    }{\eta m_*}
    +
    \frac{2\pi B_\ell K_{\mathrm{aug}}}{m_*}
        \sqrt{\bar h(s,t)},
\end{aligned}
\]
and
\[
    \|\tilde r^t-\tilde r^s\|_{L^2(\tilde{\mathbb P})}
    \leq
    K_{\mathrm{sc}}
    \|\tilde\theta^t-\tilde\theta^s\|_{\mathcal H}.
\]
Both asserted estimates therefore hold with
\[
    A:=
    \max\left\{
        \frac1{\eta m_*},
        \frac{2\pi B_\ell K_{\mathrm{aug}}}{m_*},
        K_{\mathrm{sc}}
    \right\}.
\]
All constants are deterministic and independent of $s$ and $t$.

\end{itemize}

\end{proof}

\begin{proof}[Proof of Theorem~\ref{thm:effective-process-convergence}]
By \eqref{eq:mean-field-increment-summability},
\[
    \sup_{k\geq T}
        \|\tilde\theta^{k+1}-\tilde\theta^k\|_{\mathcal H}
    \longrightarrow0.
\]
Together with the configuration estimate
\eqref{eq:effective-configuration-cauchy},
the two-time stability bound
\eqref{eq:effective-two-time-stability-main} gives
\[
\begin{aligned}
    \sup_{s,t\geq T}
        \|\tilde\theta^t-\tilde\theta^s\|_{\mathcal H}
    \leq
    A\left(
        2\sup_{k\geq T}
            \|\tilde\theta^{k+1}-\tilde\theta^k\|_{\mathcal H}
        +\sqrt{\sup_{s,t\geq T}\bar h(s,t)}
    \right)
    \longrightarrow0.
\end{aligned}
\]
Thus $(\tilde\theta^t)$ is Cauchy in $\mathcal H$.
Since $\mathcal H$ is complete, there exists
\[
    \tilde\theta^\infty
    =
    (\bar b^\infty,\tilde w^\infty)\in\mathcal H
\]
such that
$\|\tilde\theta^t-\tilde\theta^\infty\|_{\mathcal H}\to0$.
Similarly, \eqref{eq:effective-score-stability-main} gives
\[
    \sup_{s,t\geq T}
        \|\tilde r^t-\tilde r^s\|_
            {L^2(\tilde{\mathbb P})}
    \leq
    A\sup_{s,t\geq T}
        \|\tilde\theta^t-\tilde\theta^s\|_{\mathcal H}
    \longrightarrow0.
\]
Completeness of $L^2(\tilde{\mathbb P})$ therefore
gives a limit
$\tilde r^\infty\in L^2(\tilde{\mathbb P})$.
Finally, \eqref{eq:effective-switching-summability} gives
\[
\begin{aligned}
    \sum_{t=1}^{\infty}
    \tilde{\mathbb P}\left(
        \tilde\Delta=0,\,
        \tilde{\mathfrak c}^{\,t+1}
        \neq\tilde{\mathfrak c}^{\,t}
    \right)
    =
    (1-\bar\rho)\sum_{t=1}^{\infty}\bar h(t,t+1)
    <\infty.
\end{aligned}
\]
By the first Borel--Cantelli lemma, almost surely only finitely
many configuration changes occur on $\{\tilde\Delta=0\}$.
The configuration is therefore eventually constant on this event.

Define
\[
    \tilde{\mathfrak c}^{\,\infty}
    :=
    \begin{cases}
        \displaystyle
        \limsup_{t\to\infty}\tilde{\mathfrak c}^{\,t},
        &\tilde\Delta=0,\\[6pt]
        \mathbf1_{\{\tilde r^\infty\geq\kappa_+\}}
        -
        \mathbf1_{\{\tilde r^\infty\leq\kappa_-\}},
        &\tilde\Delta=1.
    \end{cases}
\]
This is a measurable $\{-1,0,+1\}$-valued random variable.
On $\{\tilde\Delta=0\}$, it agrees almost surely with
$\tilde{\mathfrak c}^{\,t}$ for all sufficiently large $t$.
Its value on $\{\tilde\Delta=1\}$ specifies an extension
to the whole sample-side probability space.
\end{proof}

\begin{proof}[Proof of Corollary~\ref{cor:mean-field-selection-stabilization}]
Since $\tilde{\mathbb P}(\tilde\Delta=0)=1-\bar\rho>0$, we have
\[
    \bar\omega^t
    =
    \frac{
        \tilde{\mathbb E}\left[
            (1-\tilde\Delta)
            \mathbf1_{\{\tilde{\mathfrak c}^{\,t}\neq0\}}
        \right]
    }{1-\bar\rho},
    \qquad
    \bar e^t
    =
    \frac{
        \tilde{\mathbb E}\left[
            (1-\tilde\Delta)
            \mathbf1_{\{\tilde{\mathfrak c}^{\,t}=-\tilde Y\}}
        \right]
    }{1-\bar\rho}.
\]
The masked indicators converge almost surely to their
counterparts with $\tilde{\mathfrak c}^{\,t}$ replaced by
$\tilde{\mathfrak c}^{\,\infty}$: on $\{\tilde\Delta=0\}$,
this follows from eventual stabilization, while on
$\{\tilde\Delta=1\}$ both indicators vanish.
Bounded convergence therefore gives the asserted limits
of $\bar\omega^t$ and $\bar e^t$.

Each indicator can change only when the configuration changes.
Consequently,
\[
\begin{aligned}
    |\bar\omega^{t+1}-\bar\omega^t|
      +|\bar e^{t+1}-\bar e^t|
    \leq
    \frac{2}{1-\bar\rho}
    \tilde{\mathbb E}\left[
        (1-\tilde\Delta)
        \mathbf1_{\{
            \tilde{\mathfrak c}^{\,t+1}
            \neq\tilde{\mathfrak c}^{\,t}
        \}}
    \right]
    =2\bar h(t,t+1).
\end{aligned}
\]
Summing over $t$ and applying
\eqref{eq:effective-switching-summability}
proves finite total variation of both sequences.
Suppose additionally that
\[
    \tilde{\mathbb P}
        (\tilde r^\infty\in\{\kappa_-,\kappa_+\})=0.
\]
The $L^2(\tilde{\mathbb P})$ convergence of the scores
admits a deterministic subsequence $t_k\to\infty$ such that
$\tilde r^{t_k}\to\tilde r^\infty$ almost surely.
The threshold map is continuous away from
$\{\kappa_-,\kappa_+\}$, so
\[
    \tilde{\mathfrak c}^{\,t_k}
    \longrightarrow
    \mathbf1_{\{\tilde r^\infty\geq\kappa_+\}}
    -
    \mathbf1_{\{\tilde r^\infty\leq\kappa_-\}}
    \qquad
    \tilde{\mathbb P}\text{-almost surely}.
\]
On $\{\tilde\Delta=0\}$, eventual stabilization also gives
$\tilde{\mathfrak c}^{\,t_k}
\to\tilde{\mathfrak c}^{\,\infty}$ almost surely.
Comparing these limits proves the desired identity on
$\{\tilde\Delta=0\}$.
On $\{\tilde\Delta=1\}$, the identity holds by the chosen
extension of $\tilde{\mathfrak c}^{\,\infty}$.
It therefore holds $\tilde{\mathbb P}$-almost surely.
\end{proof}
